%% ch13 --- Pendulums and planets
%%
%% Written September 2026. Every number the reader cannot do in their head
%% comes from code/measure/ch13.py, which imports the SAME functions the
%% listings print (code/ch13/). The double pendulum calls sin and cos, whose
%% last binary digit can differ between machines, so nothing is committed
%% from late in a chaotic run: the parting time is taken long before a
%% last-digit difference could matter, and the script refuses to write any
%% printed number that a changed step or a release moved by 1e-12 rad would
%% change.
%%
%% WHAT THIS CHAPTER MAY NOT SAY. The Lyapunov exponent and the horizon
%% formula are Chapter 7's; they are used here, not derived. The strange
%% attractor and the idea that statistics survive when paths do not are
%% Chapter 9's; this chapter only points at them. Shadowing -- whether a
%% computed chaotic path is near SOME true one -- is Chapter 16's, and the
%% note in section 3 stops short of it. The three-body problem is told as
%% history, not simulated.
%%
%% Twin: chapters/pl/ch13-mechanics.tex. Section for section, macro for
%% macro, number for number; tools/parity.py compares them.

\chapter{Pendulums and planets}
\label{ch:mechanics}
\index{double pendulum}\index{solar system}

\noindent
A grandfather clock keeps time because a pendulum is the most regular thing
most people will ever own. Pull the bob aside, let go, and it swings to and
fro at the same beat for as long as the clock is wound.
Chapter~\ref{ch:motion} showed why: the state of a pendulum is two numbers,
and without friction they go round the same closed loop for ever. The
clockwork universe of Chapter~\ref{ch:models} was built on exactly this kind
of mechanics \dash{} Newton's laws, a comet returning on schedule.

Now hang a second pendulum from the bottom of the first, with a hinge where
the two meet, and let go of it from high up. The rule is still Newton's, it
is still exact, and nothing in it is random. What does it do? This chapter
answers with a measurement, and then follows the same question from a
pendulum on a desk to the planets round the Sun, where it was first asked.

\begin{outcomes}
\outcome{step a double pendulum with RK4 and read its state from a printout}
\outcome{check a simulation of a system without friction by watching its
  energy, and say what that check does and does not prove}
\outcome{measure how long two releases a small distance apart take to part,
  and show that the answer does not depend on the size of the step}
\outcome{tell a gentle release from a wild one, and explain how a system that
  keeps its energy can still be chaotic}
\outcome{separate what has been measured about the solar system's chaos from
  what is folklore}
\end{outcomes}

\section{A pendulum on a pendulum}
\label{sec:ch13-double}
\index{double pendulum!state}

\noindent
Take two rods a metre long and hinge the second to the end of the first. Put
a one-kilogram bob at the hinge and another at the far end; the bobs carry
all the weight and the rods count as weightless. Hold both arms up,
120° from hanging straight down, and let go.

\begin{think}
One pendulum swings with perfect regularity. Do you expect two, one hanging
from the other, to settle into some pattern that repeats \dash{} more
complicated than one pendulum's, but regular all the same?
\end{think}
\reveal{No. Let go from high up, the double pendulum does not settle into any
repeating pattern you can find: the lower arm whips round, sometimes right
over the top, and the motion goes on changing. Nothing random is involved.}
The difference is the number of numbers. One pendulum's state is two numbers,
its angle and how fast the angle is changing. A double pendulum's state is
four: two angles and two rates of change. Chapter~\ref{ch:motion} gave the
reason that matters: a smooth rule whose state is two numbers cannot be
chaotic, and it takes at least three. The double pendulum has four.

Newton's laws turn into a rule of change for those four numbers, exactly as
they did for one pendulum in Chapter~\ref{ch:motion}: given the state now,
the rule says how fast each number is changing. For two arms the rule is
about fifteen lines of arithmetic with sines and cosines, and you do not need
to read it to use it. It is \api{double\_pendulum} in the book's library,
\pkg{chaoslab}, and Chapter~\ref{ch:motion}'s RK4 steps it a thousandth of a
second at a time:

\pyregion{code/ch13/release.py}{setup}{The release: four numbers, and the size of one step}{lst:ch13-setup}

\noindent
\code{math.radians} turns degrees into radians, the unit Python's sine and
cosine expect. A full turn is $2\pi$ radians, a little over six, so one
radian is about 57°. The listing then steps the state for ten seconds and
prints it every two:

\pyregion{code/ch13/release.py}{run}{Ten seconds of a double pendulum, printed every two}{lst:ch13-run}

\transcript{ch13-release}

\noindent
Read the arms first. A negative angle is a swing to the other side, and an
angle beyond 360° means the arm has gone right round: 758° at ten seconds is
two whole turns and 38° more. The lower arm is whirling. Now read the last
column.

\section{The number that must not move}
\label{sec:ch13-energy}
\index{energy}\index{simulation!energy check}

\noindent
A pendulum with no friction cannot gain or lose energy: nothing pushes it and
nothing slows it down. Its energy has two parts. The first is energy of
height: each bob contributes its weight times its height, $g \times h$ for a
kilogram, where $g$ is \num{9.81} and the height $h$ is measured upwards from
the pivot. The second is energy of motion: half of each bob's speed squared.
As the pendulum falls, height turns into speed; as it rises, speed turns back
into height. The total stays exactly the same.

\begin{think}
Suppose you let go of the pendulum with both arms horizontal, pointing the
same way, at 90°. What is its energy, with heights measured from the pivot?
\end{think}
\reveal{Zero. Both bobs are level with the pivot, so both heights are zero,
and both are still, so there is no energy of motion either.}
Heights below the pivot count as negative, so a pendulum hanging at rest has
negative energy. That does not matter: only changes in energy mean anything,
so where you measure from is your choice. For the release in the printout, a
little geometry puts the upper bob half a metre above the pivot and the lower
bob a whole metre above it, so the energy is $\num{9.81} \times (\num{0.5} +
1)$, about \val{ch13.energy.start} joules, which is where the energy column
starts. Now read down that column. Over ten seconds of whirling it changes by
a single unit in its eighth decimal place, a hundred-millionth of a joule.

That is the check, and it is the most useful habit in simulating mechanics. A
computer does not solve the rule; it takes small steps, and every step is
slightly wrong. The true pendulum keeps its energy exactly, so any change in
the simulated energy is the stepper's error, measured. A simulation of a
friction-free system whose energy drifts is not to be trusted; one whose
energy holds has passed the first test.

\begin{lab}{l13_01_energy}{The number that must not move}
Write the double pendulum's energy yourself: \code{heights} gives the height
of each bob above the pivot, and \code{energy} gives the total. The starter
explains the one step that needs care: the lower bob is carried by the upper
one, so its speed combines two swings. The test compares your energy with the
library's, and then watches it along a swing in which it must not move.
\end{lab}

\section{A millionth of a radian}
\label{sec:ch13-apart}
\index{sensitive dependence!double pendulum}

\noindent
Now let go of two double pendulums. The second is the first exactly, except
that its upper arm starts one millionth of a radian further round: at the end
of a one-metre arm, a thousandth of a millimetre, far thinner than a hair.
The two obey the same exact rule, with the same step.

\begin{think}
Both pendulums keep their energy, as the check showed, and nothing random
touches either. Will the two still swing together at the end of the ten
seconds printed above? After a minute? Or will they part, and if so, roughly
when?
\end{think}
\reveal{They part. After \val{ch13.apart.seconds} seconds they differ by more
than a tenth of a radian, about six degrees, which anyone watching would
see. A few seconds later they are doing completely different things.}

\pyregion{code/ch13/two_releases.py}{gap}{How far apart two pendulums are}{lst:ch13-gap}

\pyregion{code/ch13/two_releases.py}{run}{Two releases a millionth of a radian apart}{lst:ch13-pair}

\transcript{ch13-two-releases}

\noindent
Python writes \code{7e-05} for $7 \times 10^{-5}$, seven hundred-thousandths.
Read down the gap column. It does not grow smoothly: it rises and falls as
the arms swing, because a difference in one angle turns into a difference in
the other and back. But the trend is relentless. From a millionth to a tenth
is a factor of a hundred thousand, five factors of ten, in about ten
seconds, so on average the gap is multiplied by ten every two seconds. That
is Chapter~\ref{ch:butterfly}'s butterfly effect in a machine you could build
on a workbench.

\plotfig{ch13-two-releases}{Two double pendulums let go a millionth of a
radian apart. Above: the lower arm of each, indistinguishable for about ten
seconds. Below: the gap between them on a logarithmic scale, climbing
unevenly until it is large enough to see.}{two releases parting}

\noindent
Here is the doubt that should occur to you. The computer's pendulum is not a
real one: every step is a little wrong. What if the parting is made by those
errors and not by the pendulum? There is a direct test. If the stepper made
the parting, a more accurate stepper would move it. So run the same two
releases with four step sizes, and watch the energy at the same time:

\pyregion{code/ch13/step_check.py}{parting}{The same pair, with any step size}{lst:ch13-parting}

\transcript{ch13-step-check}

\noindent
Each halving of the step divides the energy error by about
\val{ch13.energy.shrink}, as Chapter~\ref{ch:motion}'s RK4 recipe promised,
and the moment the pair parts does not move by a hundredth of a second. At
the book's step of \val{ch13.dt} seconds the energy has moved by about
\val{ch13.energy.moved} joule before the pair parts: about one part in
\val{ch13.energy.parts}. The parting is the pendulum's, not the computer's.

\plotfig{ch13-energy-error}{How far the simulated energy wanders from its
true, constant value, for four step sizes. Each halving of the step lowers
the whole curve by about the same factor, and the dotted line marks the moment
the two releases part, which is the same for all four.}{the energy check}

\begin{note}
The check has a limit worth knowing. Compute one single release twice, with
the book's step and with half of it, and the two computations part too: by a
tenth of a radian after about \val{ch13.stepper.apart} seconds. Every step's
tiny error is a nudge, and chaos grows it like any other nudge. It starts far
smaller than a millionth of a radian, which is why our pair parts first and
every step size agrees on when. After that, the computed swing is no longer
the one this pendulum would make; Chapter~\ref{ch:computers} asks what can
still be trusted then.
\end{note}

\begin{trapbox}
\textbf{\enquote{A pendulum that keeps its energy, with no friction and
nothing random in it, cannot be chaotic.}} It needs neither friction nor
chance. Keeping its energy does pin a pendulum down, but only partly. One
pendulum's state is two numbers; fixing the energy leaves it a single
direction to move in, round one closed loop for ever, and that is why it
repeats. A double pendulum's state is four numbers; fixing the energy leaves
three directions, and three is enough for chaos. The energy check proves that
the simulation is honest. It does not make the motion predictable.
\end{trapbox}

\begin{lab}{l13_02_apart}{When do two releases part?}
Turn the listing into a function. \code{seconds\_until\_apart} steps two
pendulums side by side and returns the first moment they differ by more than a
tolerance, or \code{None} if they are still together when time runs out. Then
use it: how much longer do releases a billionth of a radian apart take to
part than releases a millionth apart? Compare your answer with
Chapter~\ref{ch:lyapunov}'s arithmetic of horizons.
\end{lab}

\section{Gentle swings and wild ones}
\label{sec:ch13-islands}
\index{KAM theorem}\index{islands of regular motion}

\noindent
Everything so far was a release from high up, with a lot of energy. A double
pendulum can also be let go gently.

\begin{think}
Release the same pendulum with both arms at 10° instead of 120°: same rule,
same machine, same nudge of a millionth of a radian. Does the gap still grow
to something anyone could see within twenty seconds?
\end{think}
\reveal{No. At 10° the two stay within a few millionths of a radian of each
other: the nudge barely grows. The same machine is calm at small swings and
wild at large ones.}

\pyregion{code/ch13/gentle_wild.py}{fate}{What a pair does, released at one angle}{lst:ch13-fate}

\transcript{ch13-gentle-wild}

\noindent
Up to \val{ch13.calm.upto}°, every pair is still together after
\val{ch13.window} seconds, and even at \val{ch13.calm.upto}° the largest gap
is \val{ch13.calm.gap} radians: the nudge has grown a few times, not a
hundred thousand times. From \val{ch13.wild.from}° on, every pair parts. The
plot fills in the angles between.

\plotfig{ch13-gentle-wild}{Each dot is a release of both arms at one angle,
with a copy a millionth of a radian further round. Small releases stay
together for twenty seconds; large ones part, and the change comes over a
few degrees.}{calm and wild releases}

\noindent
This is how chaos lives in a system without friction. With little energy the
double pendulum is nearly two simple pendulums, one swing with the arms moving
together and one with them moving against each other, and the motion is
regular. Give it more energy and chaos appears; give it more still and the
chaos takes over, with regular motion surviving only in pockets. Physicists
call the picture a chaotic sea with islands of regular motion in it.

\begin{rigourbox}
This book does not prove that the double pendulum is chaotic; it measures
sensitivity at a handful of releases. The theorem behind the sea and its
islands is called KAM, after Andrey Kolmogorov, who announced it in 1954, and
Vladimir Arnold and Jürgen Moser, who proved it in the early 1960s. It says
that when a regular, friction-free system is disturbed a little, most of its
regular motions survive, slightly bent, and chaos is confined to thin layers
between them; as the disturbance grows, the layers widen. Its proof is
graduate mathematics, found in books on the mechanics of such systems.
\end{rigourbox}

\noindent
One more difference from the chaos of Chapter~\ref{ch:lorenz}. Lorenz's
equations have friction in them, and friction squeezes every start onto the
same attractor. The double pendulum has none. Nothing is squeezed, every
start keeps its own energy for ever, and there is no attractor that every
start ends up on. What there is instead is a fence: the energy says which
states can ever be reached, and the motion, however wild, stays inside it.

\section{Two bodies, three bodies}
\label{sec:ch13-three}
\index{three-body problem}\index{Poincaré, Henri}

\noindent
The clockwork of Chapter~\ref{ch:models} was the solar system, and its
triumph was the problem of two bodies: one planet and the Sun. Johannes
Kepler found early in the seventeenth century that planets move on ellipses,
and Isaac Newton showed in 1687 that his law of gravity makes two bodies do
exactly that. Two bodies have a formula: an ellipse, traced for ever. Every
question about them has an answer you can write down.

Add a third body and there is no formula of that kind at all. In 1889 Henri
Poincaré won a prize offered in the name of King Oscar II of Sweden for work
on the problem of several bodies moving under their own gravity. While his
prize memoir was being prepared for publication, an error was found in it.
Correcting it, Poincaré found that the orbits near a simple, repeating motion
of three bodies can wind about it in a tangle so intricate that he later
wrote he would not even try to draw it. The corrected memoir appeared in
1890, and what it describes is recognised today as the first glimpse of
chaos: a rule as exact as Newton's, followed by motion that no formula
follows.

In 1908, in a book written for the general reader, Poincaré put the lesson in
words this book has been measuring: a cause too small for us to notice can
decide an effect we cannot miss, and then prediction becomes impossible.

\mermaidfig{ch13-bodies}{Two bodies have a formula; three have none; the
planets are chaotic over tens of millions of years and yet stay where they
are.}{two bodies, three bodies, the planets}

\section{Is the solar system a clock?}
\label{sec:ch13-solar}
\index{Laskar, Jacques}\index{Hyperion}

\noindent
For two centuries the solar system was the model of a perfect clock. In the
late eighteenth century Joseph-Louis Lagrange and Pierre-Simon Laplace
calculated, to the accuracy their methods allowed, that the sizes of the
planets' orbits wobble but do not drift. Laplace's demon in
Chapter~\ref{ch:models} had the planets as its evidence.

In 1989 Jacques Laskar computed the slow wobbles of the planets' orbits over
two hundred million years and found the motion of the inner planets
\dash{} Mercury, Venus, Earth and Mars \dash{} chaotic. The Lyapunov time he
found, in the sense of Chapter~\ref{ch:lyapunov}, was about five million
years: an error in where the planets are is multiplied by $e$, about
\num{2.7}, every five million years.

\begin{think}
A hundred million years is twenty of those Lyapunov times. By roughly how
much is an error multiplied over twenty of them: a hundred times, a million
times, or a billion times?
\end{think}
\reveal{About \val{ch13.solar.factor} times: nearer a billion than a million.
An error of one metre in where Earth is today would grow to
\val{ch13.solar.km} kilometres, more than the distance to the Moon.}

\begin{trapbox}
\textbf{\enquote{The solar system is a perfect clock.}} Over a human lifetime,
or a few thousand years, it is as good a clock as anything we have: eclipses
are predicted centuries ahead. Over tens of millions of years it cannot be
read. No measurement of the planets today, and no computer, can say where
Earth will be on its orbit a hundred million years from now. The rule is
Newton's and it is exact; the start is not, and five million years at a time
the error in it is multiplied by $e$.
\end{trapbox}

\noindent
Unpredictable is not the same as unstable. Laskar's result is not that the
planets fly apart, and longer computations since, run over billions of years,
find the planets keeping to orbits of much the same size and shape in almost
every run. We cannot say where Earth will be on its orbit in a hundred
million years; we can say it will very probably still be on an orbit much
like this one. That is Chapter~\ref{ch:lorenz}'s lesson in a new place: the
path cannot be predicted, and the region it moves in can.

Chaos shows itself faster in some corners of the solar system. Saturn's moon
Hyperion is a lumpy body on a stretched orbit, and calculations predict that
its spin cannot settle into a steady rotation but tumbles chaotically, tugged
one way and then another by Saturn. Observations from the Voyager and
Cassini spacecraft fit that prediction: Hyperion has no steady day.

\begin{worldbox}
Geologists date layers of sediment by the climate rhythms written into them.
Slow cycles in the shape of Earth's orbit and the tilt of its axis, named
after Milutin Milanković, change how sunlight falls on the planet, and the
climate leaves matching bands in the rock. To turn bands into dates you need
Earth's orbit far into the past, and geologists take it from computations by
Laskar's group. Those computations are trusted, cycle by cycle, back to about
fifty million years and not further, because of the chaos in this section.
\end{worldbox}

\begin{summarybox}
\begin{itemize}
\item A \textbf{double pendulum}'s state is four numbers, two angles and two
  rates of change. Its rule comes from Newton's laws, is exact and has nothing
  random in it, and RK4 steps it.
\item In a system without friction the \textbf{energy} must not change, so the
  energy of a simulation measures the stepper's error. At the book's step it
  moves by about one part in \val{ch13.energy.parts}.
\item Two releases a millionth of a radian apart part after
  \val{ch13.apart.seconds} seconds, and the moment does not change when the
  step is halved: the \textbf{parting is physics}, not numerical error.
  Keeping energy does not prevent chaos; fixing the energy of a four-number
  state still leaves three directions to move in.
\item Gentle releases are \textbf{regular} and wild ones chaotic: a chaotic
  sea with islands of regular motion, the picture the KAM theorem describes.
  With no friction there is no attractor; the energy fences the motion in.
\item Two bodies have a formula, the ellipse; three do not.
  \textbf{Poincaré}'s corrected prize memoir of 1890 was the first glimpse of
  chaos.
\item Laskar found the inner planets chaotic with a \textbf{Lyapunov time} of
  about five million years, so their positions cannot be predicted a hundred
  million years ahead. Chaotic is not unstable: their orbits stay much as they
  are.
\end{itemize}
\end{summarybox}

\canyou

\begin{checkpoint}
\item Why can a double pendulum be chaotic when a single pendulum cannot, even
  though neither loses any energy?
  \answerto{A single pendulum's state is two numbers; with the energy fixed it
  has one direction left and must go round the same loop for ever. A double
  pendulum's state is four numbers; with the energy fixed it still has three,
  which is enough for chaos.}
\item A double pendulum is let go from rest with both arms pointing straight
  up. What is its energy, measuring heights from the pivot?
  \answerto{The bobs are one and two metres above the pivot, so the energy is
  $\num{9.81} \times (1 + 2)$, about \num{29.4} joules.}
\item Your simulation of a friction-free machine loses a tenth of its energy
  in a minute. What do you conclude, and what do you do?
  \answerto{The stepper's error is far too large to trust the result. Make the
  step smaller until the energy holds, then check that the answer you care
  about does not change when the step is halved again.}
\item Two releases part after ten seconds. You halve the step and they now
  part after three. What does that tell you?
  \answerto{That the parting was made, at least partly, by the stepper's
  errors rather than by the pendulum, so neither number can be trusted. A
  parting that is physics does not move when the step changes.}
\item The double pendulum has no attractor. What keeps its motion from going
  anywhere at all?
  \answerto{Its energy. With no friction the energy never changes, and it
  decides which states can ever be reached; the motion stays inside that set.}
\item Laskar's Lyapunov time for the inner planets is about five million
  years. By roughly what factor is an error multiplied in ten million years?
  \answerto{By $e$ twice: $\num{2.7} \times \num{2.7}$, a little over seven.}
\item Does the chaos of the solar system mean the planets will fly apart?
  \answerto{No. It means their positions cannot be predicted beyond some tens
  of millions of years. Computations over billions of years find their orbits
  keeping much the same size and shape in almost every run.}
\end{checkpoint}
